% !TeX root = ../../Book3.tex
% 纲要标题：### 1.4 环同态与素谱
% 纲要位置：工作记录/计划纲要.md，第 1834 行。
% 纲要细目（写作范围）：
% * 素理想收缩与连续映射
% * 商环对应闭子空间
% * 主开集 \(D(f)\) 对应 \(R_f\)；一般局部化的谱对应与乘法集不交的素理想子空间，未必是开子空间
% ---

\section{环同态与素谱}
\label{b3:sec:0104}

环同态将元素沿一个方向送出，素理想则沿收缩反向移动。商环与局部化是理解这个反向关系的两种基本操作：前者给出闭子空间，后者保留避开指定分母的点。关于元素 \(f\) 的各次幂的局部化称为单元素局部化，它对应基本开集 \(D(f)\)；一般乘法集给出的子空间却未必是开集。

\subsection{素理想收缩与谱上的连续映射}
\label{b3:subsec:0104:01}

\begin{samepage}
\begin{proposition}\label{b3:prop:spectrum-functor}
设 \(R,S\) 为交换含幺环，\(\varphi:R\rightarrow S\) 为保单位的环同态。它给出映射
\[
\varphi^*:\operatorname{Spec}S\longrightarrow\operatorname{Spec}R,
\qquad\mathfrak q\longmapsto\varphi^{-1}(\mathfrak q).
\]
它连续。对每个理想 \(I\subseteq R\) 和每个元素 \(f\in R\)，有
\begin{equation}\label{b3:eq:spectrum-pullback}
(\varphi^*)^{-1}\bigl(V(I)\bigr)=V(IS),\qquad
(\varphi^*)^{-1}\bigl(D(f)\bigr)=D\bigl(\varphi(f)\bigr),
\end{equation}
其中 \(IS\) 是由全部 \(\varphi(a)\)、\(a\in I\)，在 \(S\) 中生成的理想。恒等同态诱导恒等映射。若 \(A\) 也是交换含幺环，\(\psi:S\rightarrow A\) 为保单位的环同态，则 \((\psi\circ\varphi)^*=\varphi^*\circ\psi^*\)。
\end{proposition}
\end{samepage}
\begin{proof}
因 \(1\) 不在 \(\mathfrak q\) 中，收缩是真理想；若 \(ab\) 的像在 \(\mathfrak q\)，素性使 \(\varphi(a)\) 或 \(\varphi(b)\) 在其中，所以收缩为素理想。包含 \(I\subseteq\varphi^{-1}(\mathfrak q)\) 等价于 \(IS\subseteq\mathfrak q\)，给出第一式和连续性。对主理想取补集得到第二式。最后两项由原像的恒等及复合规则直接得到。
\end{proof}

极大理想的收缩未必极大。例如 \(\mathbb Z\hookrightarrow\mathbb Q\) 把域的唯一素理想 \((0)\) 收缩为 \(\mathbb Z\) 的零理想；后者不是极大理想。这也是素谱采用全部素理想而不只采用极大理想的一个理由。

\subsection{商环与闭子空间}
\label{b3:subsec:0104:02}

\begin{theorem}\label{b3:thm:spec-quotient}
设 \(R\) 为交换含幺环，\(I\) 为 \(R\) 的理想。商映射 \(R\rightarrow R/I\) 诱导同胚
\[
\operatorname{Spec}(R/I)\longrightarrow V(I)\subseteq\operatorname{Spec}R,
\qquad\overline{\mathfrak p}\longmapsto\mathfrak p.
\]
这里 \(\mathfrak p\) 是 \(\overline{\mathfrak p}\) 的原像，逆映射为 \(\mathfrak p\mapsto\mathfrak p/I\)。
\end{theorem}
\begin{proof}
理想的商对应和商整环判据给出所列互逆双射。对包含 \(I\) 的理想 \(J\)，\(R/I\) 的闭集 \(V(J/I)\) 对应为 \(V(J)\)，它已包含于 \(V(I)\)；另一方面，\(V(I)\) 中任一子空间闭集可写为 \(V(I)\cap V(K)=V(I+K)\)，也如此得到。故双射及其逆都连续。
\end{proof}

取 \(I=\operatorname{Nil}(R)\)，所有素理想都包含它，因而得到 \(\operatorname{Spec}R_{\mathrm{red}}\cong\operatorname{Spec}R\)。底拓扑空间看不见被取商消去的幂零元。这不表示两个环相同；后面讨论加厚点与纤维时，会反复用到这一区别。

\subsection{主开集与一般局部化的谱}
\label{b3:subsec:0104:03}

设 \(R\) 为交换含幺环，\(T\subseteq R\) 为含一的乘法集，允许 \(0\in T\)，此时局部化为零环。第一卷定理13.2.2构造了环 \(T^{-1}R\)，命题13.2.3给出零判据：\(a/t=0\) 当且仅当存在 \(u\in T\) 使 \(ua=0\)。下面把第一卷定理13.4.3的局部化素理想对应补上拓扑说明。

\begin{theorem}\label{b3:thm:spec-localization}
设 \(R\) 为交换含幺环，\(T\subseteq R\) 为含一的乘法集。局部化映射 \(\iota:R\rightarrow T^{-1}R\) 诱导同胚
\[
\operatorname{Spec}(T^{-1}R)\longrightarrow
X_T=\{\,\mathfrak p\in\operatorname{Spec}R:\mathfrak p\cap T=\varnothing\,\},
\]
右侧赋予子空间拓扑。对应的逆为 \(\mathfrak p\mapsto T^{-1}\mathfrak p\)。特别地，对每个 \(f\in R\)，令 \(R_f=\{1,f,f^2,\ldots\}^{-1}R\)，则 \(\operatorname{Spec}R_f\cong D(f)\)。
\end{theorem}
\begin{proof}
局部化中的素理想不能包含可逆分母的像，故其收缩与 \(T\) 不交。若 \(\mathfrak p\cap T=\varnothing\)，记 \(\overline T\) 为 \(T\) 在 \(R/\mathfrak p\) 中的像。由第一卷定理13.2.4的局部化泛性质，商映射诱导满同态
\[
\theta:T^{-1}R\longrightarrow\overline T^{-1}(R/\mathfrak p),
\qquad a/t\longmapsto\overline a/\overline t.
\]
满射性来自分子、分母都能取原像。局部化的零判据说明，\(\theta(a/t)=0\) 当且仅当某个 \(u\in T\) 使 \(ua\in\mathfrak p\)；此时 \(a/t=(ua)/(ut)\in T^{-1}\mathfrak p\)。反过来，\(\theta\) 把生成扩张理想的每个 \(b/1\)、\(b\in\mathfrak p\)，都送到零。因此 \(\ker\theta=T^{-1}\mathfrak p\)，由环的第一同构定理得到
\[
(T^{-1}R)/(T^{-1}\mathfrak p)\cong\overline T^{-1}(R/\mathfrak p).
\]
右端是非零整环，因为 \(R/\mathfrak p\) 是整环且 \(\overline T\) 不含零。所以 \(T^{-1}\mathfrak p\) 为素理想。前面的核计算还说明，若 \(a/1\in T^{-1}\mathfrak p\)，则某个 \(ua\in\mathfrak p\)，其中 \(u\in T\)；因 \(u\notin\mathfrak p\)，素性迫使 \(a\in\mathfrak p\)。因此扩张后的收缩恢复 \(\mathfrak p\)。

反过来，设 \(\mathfrak q\) 为 \(T^{-1}R\) 的理想，并令 \(J=\iota^{-1}(\mathfrak q)\)。若 \(a/t\in\mathfrak q\)，则 \(a/1=(t/1)(a/t)\in\mathfrak q\)，所以 \(a\in J\)，从而 \(a/t\in T^{-1}J\)。若 \(a\in J\)、\(t\in T\)，则 \(a/t=(1/t)(a/1)\in\mathfrak q\)。故 \(T^{-1}J=\mathfrak q\)，证明了收缩后扩张也恢复原理想。这验证了素理想之间的双射。

由\cref{b3:prop:spectrum-functor}，该映射连续。局部化中的基本开集 \(D(a/t)\) 在上述双射下对应 \(D(a)\cap X_T\)：素理想是否含 \(a/t\)，只取决于它是否含分子。这些正是子空间的一组开集基，故逆映射也连续。取 \(T=\{1,f,f^2,\ldots\}\) 时，素理想与 \(T\) 不交恰好等价于不含 \(f\)，得到主开集的断言。
\end{proof}

一般的 \(X_T=\bigcap_{t\in T}D(t)\) 未必开。取 \(R=\mathbb Z\)、\(T=\mathbb Z\setminus\{0\}\)，其像只有一般点 \((0)\)。任何包含它的基本开集 \(D(n)\)、\(n\ne0\)，都还包含一个 \((p)\)，其中素数 \(p\) 不整除 \(n\)；这样的素数存在，因为有限多个素数不能穷尽全部素数。因此单点 \(\bigl\{(0)\bigr\}\) 不是开集，尽管它与 \(\operatorname{Spec}\mathbb Q\) 同胚。

\ProseParagraphBreak
\cref{b3:tab:spectrum-operations}将上述操作并列。表中 \(R\) 为交换含幺环，\(I\) 为理想，\(f\in R\)，\(\mathfrak p\) 为素理想，\(T\) 为含一的乘法集；各行都经收缩把新环的谱视为原谱的子空间。在 \(\mathfrak p\) 处局部化保留包含于 \(\mathfrak p\) 的素理想，取商 \(R/\mathfrak p\) 则保留包含 \(\mathfrak p\) 的素理想，两个方向相反。

\begin{table}[htbp]
\centering
\caption{取商与局部化在素谱中保留的点}
\label{b3:tab:spectrum-operations}
\begin{tabularx}{\textwidth}{@{}l>{\raggedright\arraybackslash}p{.29\textwidth}>{\raggedright\arraybackslash}X@{}}
\toprule
环的操作 & 原环中保留的素理想 \(\mathfrak q\) & 在 \(\operatorname{Spec}R\) 中的像 \\
\midrule
\(R\rightarrow R/I\) & \(I\subseteq\mathfrak q\) & \(V(I)\)，闭子空间 \\
\(R\rightarrow R_f\) & \(f\notin\mathfrak q\) & \(D(f)\)，基本开集 \\
\(R\rightarrow R_{\mathfrak p}\) & \(\mathfrak q\subseteq\mathfrak p\) & \(\mathfrak p\) 及其一般化组成的子空间，未必开 \\
\(R\rightarrow T^{-1}R\) & \(\mathfrak q\cap T=\varnothing\) & \(\bigcap_{t\in T}D(t)\)，未必开 \\
\bottomrule
\end{tabularx}
\end{table}

第二章将把这些环上的分式构造推广到模。
